\documentclass{article}
\usepackage{amsfonts}
\usepackage{bm}
\usepackage{algorithm}
\usepackage{algorithmic}
\usepackage{mathtools}
\usepackage{graphicx}
\usepackage{booktabs}
\usepackage{xcolor}
\usepackage{hyperref}
\usepackage{amsmath, amssymb, amsthm}
\usepackage{geometry}
\usepackage{natbib}
\geometry{margin=1in}

\newtheorem{theorem}{Theorem}
\newtheorem{proposition}[theorem]{Proposition}
\newtheorem{corollary}[theorem]{Corollary}
\newtheorem{lemma}[theorem]{Lemma}
\newtheorem{remark}{Remark}
\newtheorem{definition}{Definition}
\newtheorem{assumption}{Assumption}

\DeclareMathOperator{\tr}{tr}
\DeclareMathOperator{\closure}{closure}
\DeclareMathOperator{\ilr}{ilr}
\DeclareMathOperator{\diag}{diag}
\DeclareMathOperator{\softmax}{softmax}
\DeclareMathOperator{\vect}{vec}

\begin{document}

\title{Structural Topic Modeling in Aitchison Geometry:\\
  An Analytical Spectral Solution via ILR-Spectral-GSCA}
\date{\today}
\maketitle

\begin{abstract}
We introduce a structural topic model that operates natively in the
Aitchison geometry of the simplex. Document--topic proportions are
represented in isometric log-ratio (ILR) coordinates, and covariates are
linked to topic prevalence through a Generalized Structured Component
Analysis (GSCA) formulation. By profiling out all nuisance parameters in
closed form, we reduce the estimation problem to a single eigenvalue
decomposition of a covariate-augmented document similarity matrix.
The resulting estimator is non-iterative, globally optimal for the
profiled linearised objective, and yields topic proportions on the
probability simplex by construction---not through a post-hoc
transformation, but as the natural inverse of the coordinate system in
which estimation is performed. We establish global optimality for the
linearised ILR-Spectral-GSCA objective, clarify the identifiability conditions
of the spectral solution under both simple and degenerate spectra, and
derive consistency and asymptotic normality of the covariate path
coefficients with an explicit analytical covariance formula based on
eigenvector perturbation theory. A linearisation error bound under an
eigenvector delocalisation condition quantifies the approximation quality
relative to the original simplex problem. The approach bridges three
previously disjoint frameworks---topic modeling, structural equation
modeling, and compositional data analysis---through the unifying lens of
Aitchison geometry.
\end{abstract}

% ====================================================================
\section{Introduction}
\label{sec:introduction}
% ====================================================================

[To be written.]

\input{methodology}

% ====================================================================
\section{Algorithms and Implementation}
\label{sec:algorithm}
% ====================================================================

The estimation procedure consists of two stages.
Algorithm~\ref{alg:ilr_egsca} computes the globally optimal spectral
solution for the profiled linearised objective; its output is
deterministic and non-iterative.
Algorithm~\ref{alg:refine} is an optional single-pass refinement that
re-estimates the topic--term distributions outside the ILR subspace,
improving topic exclusivity without modifying the document--topic
proportions or path coefficients.

% ---- Algorithm 1 --------------------------------------------------
\begin{algorithm}[h]
\caption{ILR-EGSCA Spectral Fit}
\label{alg:ilr_egsca}
\begin{algorithmic}[1]
\REQUIRE Document--term matrix
  $\mathbf{W}\in\mathbb{R}^{M\times N}_{\geq 0}$,
  column-centred covariate matrix
  $\mathbf{C}\in\mathbb{R}^{M\times P}$,
  number of topics $K$,
  regularisation $\lambda>0$,
  ILR contrast matrix
  $\mathbf{V}\in\mathbb{R}^{K\times(K-1)}$.
\ENSURE Topic proportions $\hat{\boldsymbol{\Pi}}\in\mathbb{R}^{M\times K}$,
  ILR loadings $\hat{\boldsymbol{\Psi}}\in\mathbb{R}^{(K-1)\times N}$,
  ILR path coefficients $\hat{\mathbf{B}}_z\in\mathbb{R}^{P\times(K-1)}$.

\medskip
\STATE \textbf{Centre.}\;
  $\widetilde{\mathbf{W}} \leftarrow \mathbf{W}
    - M^{-1}\mathbf{1}_M\mathbf{1}_M^{\!\top}\mathbf{W}$
  \hfill$\triangleright$ Prop.~\ref{prop:mu_profile}

\STATE \textbf{Truncated SVD.}\;
  $\mathbf{U}\boldsymbol{\Sigma}\mathbf{V}_W^{\!\top}
    \leftarrow \text{SVD}_r(\widetilde{\mathbf{W}})$
  \hfill$\triangleright$ randomised, rank $r$

\STATE \textbf{Covariate QR.}\;
  $\mathbf{Q}_C\mathbf{R}_C \leftarrow \text{QR}(\mathbf{C})$

\STATE \textbf{Augmented factor.}\;
  $\mathbf{H} \leftarrow [\,\mathbf{U}\boldsymbol{\Sigma}
    \;\;\;\sqrt{\lambda}\,\mathbf{Q}_C\,]
    \in\mathbb{R}^{M\times(r+P)}$
  \hfill$\triangleright$ Eq.~\eqref{eq:H_matrix}

\STATE \textbf{Reduced eigendecomposition.}\;
  $\mathbf{H}^{\!\top}\mathbf{H}
    = \mathbf{E}\,\diag(\sigma_1,\dots,\sigma_{r+P})\,
      \mathbf{E}^{\!\top}$
  \hfill$\triangleright$ $(r{+}P)\times(r{+}P)$

\STATE \textbf{ILR topic scores.}\;
  $\mathbf{Z}^* \leftarrow \mathbf{H}\mathbf{E}_{1:K-1}\,
    \diag(\sigma_1^{-1/2},\dots,\sigma_{K-1}^{-1/2})$
  \hfill$\triangleright$ Thm.~\ref{thm:spectral}

\STATE \textbf{ILR loadings.}\;
  $\hat{\boldsymbol{\Psi}} \leftarrow
    K\,\mathbf{Z}^{*\top}\widetilde{\mathbf{W}}$
  \hfill$\triangleright$ Prop.~\ref{prop:Psi_profile}

\STATE \textbf{Varimax rotation} (optional)\textbf{.}\;
  $\mathbf{R}^* \leftarrow \mathrm{varimax}
    (\hat{\boldsymbol{\Psi}}^{\!\top})$;\;
  $\mathbf{Z}^* \leftarrow \mathbf{Z}^*\mathbf{R}^*$;\;
  $\hat{\boldsymbol{\Psi}} \leftarrow
    \mathbf{R}^{*\top}\hat{\boldsymbol{\Psi}}$
  \hfill$\triangleright$ Prop.~\ref{prop:rotation_invariance}

\STATE \textbf{Path coefficients.}\;
  $\hat{\mathbf{B}}_z \leftarrow
    (\mathbf{C}^{\!\top}\mathbf{C})^{-1}
    \mathbf{C}^{\!\top}\mathbf{Z}^*$
  \hfill$\triangleright$ Prop.~\ref{prop:Bz_profile}

\STATE \textbf{Topic proportions.}\;
  $\hat{\pi}_{ik} \leftarrow
    \exp\!\bigl((\mathbf{V}\mathbf{z}_i^*)_k\bigr)\big/
    \sum_{k'}\exp\!\bigl((\mathbf{V}\mathbf{z}_i^*)_{k'}\bigr)$
  \quad$\forall\; i,k$
  \hfill$\triangleright$ ILR inverse

\RETURN $\hat{\boldsymbol{\Pi}},\;
  \hat{\boldsymbol{\Psi}},\;
  \hat{\mathbf{B}}_z,\;
  \mathbf{Z}^*$
\end{algorithmic}
\end{algorithm}

\noindent
Algorithm~\ref{alg:ilr_egsca} is \emph{non-iterative}: it consists
of a fixed sequence of matrix decompositions and linear-algebra
operations, with no convergence loop and no random initialisation.
Its output is deterministic for any given input (up to the optional
varimax step, which is itself a deterministic function of
$\hat{\boldsymbol{\Psi}}$). The computational bottleneck is the
truncated SVD in Step~2, implemented via randomised projection
\citep{halko2011finding} at cost $O(MNr)$; all subsequent steps
operate on matrices of dimension at most
$(r{+}P)\times(r{+}P)$.


% ---- Structural limitation of Phi ---------------------------------
\subsection{Topic--Term Refinement}
\label{subsec:refinement}

The ILR-based topic--term matrix
$\hat{\boldsymbol{\Phi}} = \mathbf{V}\hat{\boldsymbol{\Psi}}
+ \mathbf{1}_K\bar{\mathbf{w}}^{\!\top}$
(Corollary~\ref{cor:recovery}) constrains the $K$~topic--term
distributions to a $(K{-}1)$-dimensional affine subspace of
$\mathbb{R}^N$ anchored at the corpus mean~$\bar{\mathbf{w}}$.
This is the algebraic cost of the closed-form solution: the same
structure that permits profiling out $\boldsymbol{\Psi}$ in
Proposition~\ref{prop:Psi_profile} forces the $K$~rows of
$\hat{\boldsymbol{\Phi}}$ to share a common baseline and
differentiate only along $K{-}1$ orthogonal directions. The result is
high topic similarity and limited exclusivity relative to methods
that estimate each topic distribution freely (e.g., STM, LDA).

To address this, we introduce a single-pass refinement that
re-estimates $\hat{\boldsymbol{\Phi}}$ outside the ILR subspace
while leaving the spectral estimates
$(\hat{\boldsymbol{\Pi}}, \hat{\mathbf{B}}_z, \mathbf{Z}^*)$
unchanged. The key observation is that the document--topic
proportions $\hat{\boldsymbol{\Pi}}$ produced by
Algorithm~\ref{alg:ilr_egsca}---which lie on the simplex and enjoy
the theoretical guarantees of
Section~\ref{sec:theory}---provide a valid soft assignment of
documents to topics, from which unconstrained topic--term
distributions can be recovered in closed form.

\paragraph{Assignment step.}
Two strategies are available for computing the document--topic
assignment weights $\mathbf{Q}\in\mathbb{R}^{M\times K}$:

\begin{enumerate}
  \item[\emph{(a)}]
    \textbf{$K$-means on ILR scores} (default).
    Apply $K$-means clustering to the scaled ILR scores
    $\sqrt{M}\,\mathbf{Z}^*$ (which have approximately unit
    column variance), producing hard assignments
    $q_{ik}\in\{0,1\}$.
    This directly exploits the spectral geometry:
    clusters in ILR space correspond to groups of documents with
    similar topic compositions.

  \item[\emph{(b)}]
    \textbf{Soft E-step}.
    Compute soft assignments via
    $q_{ik}\propto\exp\bigl(\sum_n\widetilde{w}_{in}\,
    \log\hat\phi^{(0)}_{kn}\bigr)$,
    where
    $\hat{\boldsymbol{\Phi}}^{(0)}
    = \softmax(\mathbf{V}\hat{\boldsymbol{\Psi}})$
    is the ILR-based topic--word distribution.
    This is one LDA-style E-step and tends to improve semantic
    coherence at the cost of smaller gains in diversity.
\end{enumerate}

\paragraph{M-step.}
Given $\mathbf{Q}$, the refined topic--term distributions are
\begin{equation}
  \hat{\phi}^{(1)}_{kn}
  \;=\;
  \frac{\sum_{i=1}^M q_{ik}\,\widetilde{w}_{in} + \varepsilon}
       {\sum_{n'=1}^{N}\bigl(\sum_{i=1}^M q_{ik}\,
        \widetilde{w}_{in'} + \varepsilon\bigr)},
  \label{eq:mstep_phi}
\end{equation}
where $\varepsilon>0$ is a Laplace smoothing constant.
In matrix form:
$\hat{\boldsymbol{\Phi}}^{(1)}
 = \mathrm{row\text{-}normalise}\!\bigl(
     \mathbf{Q}^{\!\top}\widetilde{\mathbf{W}} + \varepsilon\bigr)$,
which costs $O(MKN)$---a single matrix multiplication.

\paragraph{Temperature sharpening.}
An optional temperature parameter $\tau\in(0,1]$ further concentrates
mass on the highest-probability terms:
\begin{equation}
  \hat{\phi}^{(\tau)}_{kn}
  \;=\;
  \frac{\bigl(\hat{\phi}^{(1)}_{kn}\bigr)^{1/\tau}}
       {\sum_{n'}\bigl(\hat{\phi}^{(1)}_{kn'}\bigr)^{1/\tau}}.
  \label{eq:temperature}
\end{equation}
Setting $\tau=1$ recovers the M-step distribution unchanged;
$\tau<1$ sharpens the distribution towards a sparse representation.

% ---- Algorithm 2 --------------------------------------------------
\begin{algorithm}[t]
\caption{Topic--Term Refinement (post-spectral)}
\label{alg:refine}
\begin{algorithmic}[1]
\REQUIRE Spectral fit from Algorithm~\ref{alg:ilr_egsca}:
  $\hat{\boldsymbol{\Pi}}$, $\hat{\boldsymbol{\Psi}}$,
  $\mathbf{Z}^*$;\;
  document--term matrix~$\mathbf{W}$;\;
  smoothing $\varepsilon>0$;\;
  temperature $\tau\in(0,1]$;\;
  method $\in\{\text{kmeans},\,\text{em}\}$.
\ENSURE Refined topic--term matrix
  $\hat{\boldsymbol{\Phi}}^{(\tau)}
   \in\mathbb{R}^{K\times N}$.

\medskip
\IF{method $=$ kmeans}
  \STATE $\mathbf{Q} \leftarrow \text{$K$-means}
    (\sqrt{M}\,\mathbf{Z}^*,\; K\text{ clusters})$
  \hfill$\triangleright$ hard assignment
\ELSE
  \STATE $\hat{\boldsymbol{\Phi}}^{(0)} \leftarrow
    \softmax(\mathbf{V}\hat{\boldsymbol{\Psi}})$
  \STATE $q_{ik} \leftarrow
    \softmax_k\!\bigl(\widetilde{\mathbf{w}}_i^{\!\top}
    \log\hat{\boldsymbol{\Phi}}^{(0)\top}\bigr)$
  \hfill$\triangleright$ soft E-step
\ENDIF
\STATE \textbf{M-step.}\;
  $\hat{\boldsymbol{\Phi}}^{(1)} \leftarrow
    \text{row-normalise}\!\bigl(
      \mathbf{Q}^{\!\top}\widetilde{\mathbf{W}}
      + \varepsilon\bigr)$
  \hfill$\triangleright$ Eq.~\eqref{eq:mstep_phi}
\STATE \textbf{Temperature.}\;
  $\hat{\phi}^{(\tau)}_{kn} \leftarrow
    \hat{\phi}^{(1)1/\tau}_{kn}\big/
    \sum_{n'}\hat{\phi}^{(1)1/\tau}_{kn'}$
  \hfill$\triangleright$ Eq.~\eqref{eq:temperature}

\RETURN $\hat{\boldsymbol{\Phi}}^{(\tau)}$
\end{algorithmic}
\end{algorithm}

\begin{remark}[Separation of estimation and refinement]
\label{rem:separation}
Algorithm~\ref{alg:refine} modifies only the topic--term matrix
$\hat{\boldsymbol{\Phi}}$. The document--topic proportions
$\hat{\boldsymbol{\Pi}}$, the path coefficients
$\hat{\mathbf{B}}_z$, and the ILR scores $\mathbf{Z}^*$ are
\emph{unchanged}. Therefore all theoretical guarantees of
Section~\ref{sec:theory}---global optimality
(Theorem~\ref{thm:global}), consistency
(Theorem~\ref{thm:consistency}), asymptotic normality
(Theorem~\ref{thm:normality}), and the linearisation bound
(Proposition~\ref{prop:linearisation_error})---remain valid for the
structural parameters
$(\hat{\boldsymbol{\Pi}},\hat{\mathbf{B}}_z)$ after refinement.
The refinement addresses only the representation of the
topic--term distributions, which are a derived quantity not covered
by the asymptotic theory.
\end{remark}

\begin{remark}[Why a single pass suffices]
\label{rem:single_pass}
In classical EM for topic models, the E-step and M-step are iterated
until convergence. Here, the document--topic proportions are
\emph{fixed} by the spectral solution and are not updated by the
M-step, so there is no E-M loop: the refinement is a single
forward pass from the spectral $\hat{\boldsymbol{\Pi}}$
(or from the $K$-means partition of $\mathbf{Z}^*$)
to the unconstrained $\hat{\boldsymbol{\Phi}}^{(\tau)}$.
The computational cost is dominated by the matrix product
$\mathbf{Q}^{\!\top}\widetilde{\mathbf{W}}\in\mathbb{R}^{K\times N}$
at $O(MKN)$; for $K$-means, the clustering step adds $O(MK^2)$
per restart, which is negligible for moderate~$K$.
\end{remark}


\subsection{Complete Pipeline}
\label{subsec:pipeline}

The recommended workflow for applied use combines both algorithms:
%
\begin{enumerate}
  \item \textbf{Spectral fit}
    (Algorithm~\ref{alg:ilr_egsca}): produces
    $\hat{\boldsymbol{\Pi}}$, $\hat{\mathbf{B}}_z$,
    $\hat{\boldsymbol{\Psi}}$, $\mathbf{Z}^*$ with full theoretical
    guarantees. Varimax rotation is applied by default.

  \item \textbf{Topic--term refinement}
    (Algorithm~\ref{alg:refine}): re-estimates
    $\hat{\boldsymbol{\Phi}}$ via $K$-means on ILR scores followed by
    a weighted M-step with Laplace smoothing and optional temperature
    sharpening. Default: $\varepsilon=10^{-4}$, $\tau=0.5$.

  \item \textbf{Inference}: standard errors for
    $\hat{\mathbf{B}}_z$ are computed from the analytical sandwich
    formula (Theorem~\ref{thm:normality}) or, in the presence of
    small eigengaps, via nonparametric bootstrap.
\end{enumerate}
%
The entire pipeline---from raw document--term matrix to topic
proportions, path coefficients with confidence intervals, and
refined topic--term distributions---requires no iterative
optimisation and no random initialisation beyond the $K$-means
step in Algorithm~\ref{alg:refine}.

\subsection{Computational Complexity}
\label{subsec:complexity}

Table~\ref{tab:complexity} summarises the cost of each step. The
dominant cost is the randomised SVD in Algorithm~\ref{alg:ilr_egsca},
Step~2, which scales as $O(MN\log r)$ via the algorithm
of~\citet{halko2011finding}. All other steps are linear in~$M$.
For comparison, STM requires $O(MK\cdot\text{iter}_{\text{EM}})$
per variational E-step and typically 50--100 EM iterations, making
it one to two orders of magnitude slower on large corpora.

\begin{table}[h]
\centering
\caption{Computational cost per step. Here $r$~is the SVD truncation
rank, $P$~the number of covariates, $K$~the number of topics, and
$N_{\text{start}}$~the number of $K$-means restarts.}
\label{tab:complexity}
\begin{tabular}{llc}
\toprule
Step & Operation & Cost \\
\midrule
\multicolumn{3}{l}{\emph{Algorithm~\ref{alg:ilr_egsca}: Spectral fit}} \\
\; 1 & Column centering & $O(MN)$ \\
\; 2 & Randomised SVD (rank $r$) & $O(MN\log r)$ \\
\; 3 & Covariate QR & $O(MP^2)$ \\
\; 4--6 & Gram eigendecomposition & $O((r{+}P)^3 + M(r{+}P)^2)$ \\
\; 7 & ILR loadings $\hat{\boldsymbol{\Psi}}$ & $O((K{-}1)MN)$ \\
\; 8 & Varimax rotation & $O(N(K{-}1)^2)$ \\
\; 9 & Path coefficients & $O(MP(K{-}1))$ \\
\; 10 & Softmax (ILR inverse) & $O(MK)$ \\
\midrule
\multicolumn{3}{l}{\emph{Algorithm~\ref{alg:refine}: Refinement}} \\
\; 1 & $K$-means on $\mathbf{Z}^*$ & $O(MK^2 N_{\text{start}})$ \\
\; 2 & M-step & $O(MKN)$ \\
\; 3 & Temperature sharpening & $O(KN)$ \\
\bottomrule
\end{tabular}
\end{table}


\bibliographystyle{apalike}
\bibliography{bibliography}
\end{document}